% shared content, included by both main.tex and preprint.tex
Raman mapping produces a hyperspectral cube, and turning that cube into numbers
depends strongly on a chain of preprocessing choices. Two further problems
appear once a laboratory holds many maps rather than one: the corpus is
analysed one file at a time, and the analysable corpus is only the subset
somebody remembered to export by hand from the instrument's proprietary
format. We present \ramansp, a compact open framework providing typed spectral
containers, composable preprocessing protocols and an analysis layer, and
extend it in three directions. First, a direct reader for the undocumented
Horiba LabSpec\,6 container, reverse engineered here and verified bit identical
against the instrument's own text export, which removes the manual export step
and takes this corpus from \Nrawmaptext{} usable map to \Nrawmap.
Second, a cross-sample knowledge graph linking every acquisition by spectral
similarity, shared unmixing endmembers and disorder-metric proximity, with each
metric edge carrying the preprocessing protocol under which it was computed.
Third, and as the main methodological contribution, spectral 3D Gaussian
splatting: a map $V(x,y,\nu)$ is represented by a differentiable field of
anisotropic ellipsoids, each localising a domain in the two spatial axes and a
band along the wavenumber axis, fitted by first-order descent on closed-form
gradients, entirely in NumPy on a CPU. On an anonymised corpus of \Nacq{}
acquisitions the graph recovers \KGcommunities{} communities at modularity
$Q=\KGmodularity$ and flags \Ndupgroups{} groups of byte-identical
duplicates, and across \SplatNfits{} maps the field reconstructs held-out
wavenumber channels to within \SplatResidNoiseMin{} to
\SplatResidNoiseMax{} of each map's own measurement noise at
\SplatComprMin{} to \SplatComprMax$\times$ compression, that is,
essentially to the noise floor.
